% Options for packages loaded elsewhere
\PassOptionsToPackage{unicode}{hyperref}
\PassOptionsToPackage{hyphens}{url}
\documentclass[
  11pt,
  a4paper,
]{article}
\usepackage{xcolor}
\usepackage[margin=18mm]{geometry}
\usepackage{amsmath,amssymb}
\setcounter{secnumdepth}{-\maxdimen} % remove section numbering
\usepackage{iftex}
\ifPDFTeX
  \usepackage[T1]{fontenc}
  \usepackage[utf8]{inputenc}
  \usepackage{textcomp} % provide euro and other symbols
\else % if luatex or xetex
  \usepackage{unicode-math} % this also loads fontspec
  \defaultfontfeatures{Scale=MatchLowercase}
  \defaultfontfeatures[\rmfamily]{Ligatures=TeX,Scale=1}
\fi
\ifPDFTeX\else
  % xetex/luatex font selection
    \setmainfont[]{Noto Serif}
    \setsansfont[]{Noto Sans}
    \setmonofont[]{Noto Sans Mono}
  \setmathfont[]{Noto Sans Math}
\fi
% Use upquote if available, for straight quotes in verbatim environments
\IfFileExists{upquote.sty}{\usepackage{upquote}}{}
\IfFileExists{microtype.sty}{% use microtype if available
  \usepackage[]{microtype}
  \UseMicrotypeSet[protrusion]{basicmath} % disable protrusion for tt fonts
}{}
\makeatletter
\@ifundefined{KOMAClassName}{% if non-KOMA class
  \IfFileExists{parskip.sty}{%
    \usepackage{parskip}
  }{% else
    \setlength{\parindent}{0pt}
    \setlength{\parskip}{6pt plus 2pt minus 1pt}}
}{% if KOMA class
  \KOMAoptions{parskip=half}}
\makeatother
\ifLuaTeX
\usepackage[bidi=basic]{babel}
\else
\usepackage[bidi=default]{babel}
\fi
\babelprovide[main,import]{russian}
\ifPDFTeX
\else
\babelfont{rm}[]{Noto Serif}
\fi
% get rid of language-specific shorthands (see #6817):
\let\LanguageShortHands\languageshorthands
\def\languageshorthands#1{}
\setlength{\emergencystretch}{3em} % prevent overfull lines
\providecommand{\tightlist}{%
  \setlength{\itemsep}{0pt}\setlength{\parskip}{0pt}}
\usepackage{bookmark}
\IfFileExists{xurl.sty}{\usepackage{xurl}}{} % add URL line breaks if available
\urlstyle{same}
\hypersetup{
  pdftitle={Билеты по дискретным случайным процессам},
  pdflang={ru-RU},
  hidelinks,
  pdfcreator={LaTeX via pandoc}}

\title{Билеты по дискретным случайным процессам}
\author{}
\date{}

\begin{document}
\maketitle

\section{Билет 12. Упорядоченные и алгебраические структуры; гильбертовы
пространства и
проекторы}\label{ux431ux438ux43bux435ux442-12.-ux443ux43fux43eux440ux44fux434ux43eux447ux435ux43dux43dux44bux435-ux438-ux430ux43bux433ux435ux431ux440ux430ux438ux447ux435ux441ux43aux438ux435-ux441ux442ux440ux443ux43aux442ux443ux440ux44b-ux433ux438ux43bux44cux431ux435ux440ux442ux43eux432ux44b-ux43fux440ux43eux441ux442ux440ux430ux43dux441ux442ux432ux430-ux438-ux43fux440ux43eux435ux43aux442ux43eux440ux44b}

Источники: Ширяев 1, гл. II, §11; Сердобольская, лекция 6; лекции ДСП.

\subsection{1. Короткий
ответ}\label{ux43aux43eux440ux43eux442ux43aux438ux439-ux43eux442ux432ux435ux442}

В этом билете нужно связать общий язык структур с вероятностной
геометрией второго порядка.

\textbf{Упорядоченная структура} - это множество с отношением порядка.
Например, для случайных величин можно писать

\[
X\le Y \quad \text{почти наверное}.
\]

\textbf{Алгебраическая структура} - это множество с операциями. В
вероятности главные примеры: линейные пространства случайных величин,
пространства \(L^p\), алгебры событий, \(\sigma\)-алгебры.

Главная часть билета: случайные величины с конечным вторым моментом
образуют гильбертово пространство

\[
L^2(\Omega,\mathcal F,P),
\]

со скалярным произведением

\[
\langle X,Y\rangle=\mathbb E[XY]
\]

в вещественном случае и нормой

\[
\|X\|_2=(\mathbb EX^2)^{1/2}.
\]

Для центрированных величин скалярное произведение становится
ковариацией:

\[
\langle X-\mathbb EX,\ Y-\mathbb EY\rangle
=
\operatorname{Cov}(X,Y).
\]

Если \(H\) - замкнутое подпространство гильбертова пространства, то для
любого \(x\) существует единственная ортогональная проекция
\(P_Hx\in H\):

\[
x-P_Hx\perp H.
\]

То есть \(P_Hx\) - ближайший к \(x\) элемент из \(H\):

\[
\|x-P_Hx\|=\min_{h\in H}\|x-h\|.
\]

В вероятности это дает геометрический смысл условного ожидания:

\[
\mathbb E[X\mid\mathcal G]
\]

\begin{itemize}
\tightlist
\item
  ортогональная проекция \(X\in L^2\) на замкнутое подпространство
  \(\mathcal G\)-измеримых величин из \(L^2\).
\end{itemize}

\subsection{2. Полный
ответ}\label{ux43fux43eux43bux43dux44bux439-ux43eux442ux432ux435ux442}

\subsubsection{2.1. Переход к геометрии случайных
величин}\label{ux43fux435ux440ux435ux445ux43eux434-ux43a-ux433ux435ux43eux43cux435ux442ux440ux438ux438-ux441ux43bux443ux447ux430ux439ux43dux44bux445-ux432ux435ux43bux438ux447ux438ux43d}

В программе этот билет стоит как переход от меры и интеграла к линейной
геометрии случайных величин. До этого случайные величины рассматривались
как измеримые функции, у которых есть распределения и моменты. Теперь
добавляется идея: если у случайных величин есть вторые моменты, то с
ними можно работать как с векторами в гильбертовом пространстве.

\subsubsection{2.2. Упорядоченные
структуры}\label{ux443ux43fux43eux440ux44fux434ux43eux447ux435ux43dux43dux44bux435-ux441ux442ux440ux443ux43aux442ux443ux440ux44b}

Сначала нужен общий язык структур.

\textbf{Упорядоченная структура} - это множество \(M\) с отношением
\(\le\). Отношение называется частичным порядком, если для любых
\(a,b,c\in M\) выполнены:

\begin{enumerate}
\def\labelenumi{\arabic{enumi}.}
\tightlist
\item
  рефлексивность:
\end{enumerate}

\[
a\le a;
\]

\begin{enumerate}
\def\labelenumi{\arabic{enumi}.}
\setcounter{enumi}{1}
\tightlist
\item
  антисимметричность:
\end{enumerate}

\[
a\le b,\ b\le a\quad\Rightarrow\quad a=b;
\]

\begin{enumerate}
\def\labelenumi{\arabic{enumi}.}
\setcounter{enumi}{2}
\tightlist
\item
  транзитивность:
\end{enumerate}

\[
a\le b,\ b\le c\quad\Rightarrow\quad a\le c.
\]

Если любые два элемента сравнимы, порядок называется линейным, или
полным. Например, на \(\mathbb R\) обычный порядок полный. На множестве
подмножеств \(\Omega\) порядок по включению

\[
A\subset B
\]

частичный: не любые два множества сравнимы.

В вероятности важны несколько порядков:

\[
A\subset B
\]

для событий,

\[
X\le Y
\]

для случайных величин, обычно в смысле почти наверное. Строго говоря,
это частичный порядок на классах эквивалентности по равенству почти
наверное; на конкретных версиях случайных величин это только
предпорядок. Также используют порядок

\[
f\le g
\]

для измеримых функций. Этот порядок нужен в теореме Беппо Леви,
монотонности интеграла и монотонности условного ожидания.

Есть также порядок на неотрицательно определенных матрицах:

\[
A\le B
\quad\Longleftrightarrow\quad
B-A\ge0
\]

в смысле квадратичных форм:

\[
x^T(B-A)x\ge0
\]

для всех \(x\). Такой порядок используется для ковариационных матриц.

\subsubsection{2.3. Алгебраические
структуры}\label{ux430ux43bux433ux435ux431ux440ux430ux438ux447ux435ux441ux43aux438ux435-ux441ux442ux440ux443ux43aux442ux443ux440ux44b}

\textbf{Алгебраическая структура} - это множество с одной или
несколькими операциями, удовлетворяющими аксиомам. Примеры:

\begin{enumerate}
\def\labelenumi{\arabic{enumi}.}
\tightlist
\item
  полугруппа - множество с ассоциативной операцией;
\item
  моноид - полугруппа с нейтральным элементом;
\item
  группа - моноид, где у каждого элемента есть обратный;
\item
  кольцо - множество с операциями сложения и умножения;
\item
  поле - кольцо, где можно делить на ненулевые элементы;
\item
  линейное пространство - множество, где можно складывать векторы и
  умножать их на числа из поля.
\end{enumerate}

\subsubsection{2.4. Линейное пространство случайных
величин}\label{ux43bux438ux43dux435ux439ux43dux43eux435-ux43fux440ux43eux441ux442ux440ux430ux43dux441ux442ux432ux43e-ux441ux43bux443ux447ux430ux439ux43dux44bux445-ux432ux435ux43bux438ux447ux438ux43d}

В этом билете главный пример - линейное пространство случайных величин.
Если \(X\) и \(Y\) - случайные величины, а \(a,b\in\mathbb R\), то

\[
aX+bY
\]

снова случайная величина. Если \(X,Y\in L^2\), то по неравенству

\[
(aX+bY)^2\le 2a^2X^2+2b^2Y^2
\]

получаем \(aX+bY\in L^2\). Значит \(L^2\) - линейное пространство.

\subsubsection{2.5. Гильбертовы
пространства}\label{ux433ux438ux43bux44cux431ux435ux440ux442ux43eux432ux44b-ux43fux440ux43eux441ux442ux440ux430ux43dux441ux442ux432ux430}

Теперь переходим к гильбертовым пространствам. Сначала определения.

Линейное пространство \(V\) над \(\mathbb R\) называется пространством
со скалярным произведением, если задана функция

\[
\langle x,y\rangle
\]

такая, что:

\begin{enumerate}
\def\labelenumi{\arabic{enumi}.}
\tightlist
\item
  линейность по первому аргументу:
\end{enumerate}

\[
\langle ax_1+bx_2,y\rangle
=
a\langle x_1,y\rangle+b\langle x_2,y\rangle;
\]

\begin{enumerate}
\def\labelenumi{\arabic{enumi}.}
\setcounter{enumi}{1}
\tightlist
\item
  симметрия:
\end{enumerate}

\[
\langle x,y\rangle=\langle y,x\rangle;
\]

\begin{enumerate}
\def\labelenumi{\arabic{enumi}.}
\setcounter{enumi}{2}
\tightlist
\item
  положительная определенность:
\end{enumerate}

\[
\langle x,x\rangle\ge0,
\qquad
\langle x,x\rangle=0\Longleftrightarrow x=0.
\]

В комплексном случае вместо симметрии берут эрмитову симметрию:

\[
\langle x,y\rangle=\overline{\langle y,x\rangle},
\]

и одно из направлений линейно, другое сопряженно-линейно.

Скалярное произведение задает норму:

\[
\|x\|=\sqrt{\langle x,x\rangle}.
\]

Пространство со скалярным произведением называется гильбертовым, если
оно полно по этой норме: всякая фундаментальная последовательность
сходится к элементу того же пространства.

\subsubsection{\texorpdfstring{2.6. Пространство
\(L^2\)}{2.6. Пространство L\^{}2}}\label{ux43fux440ux43eux441ux442ux440ux430ux43dux441ux442ux432ux43e-l2}

Главный пример для вероятности:

\[
L^2(\Omega,\mathcal F,P)
=
\{X:\mathbb E X^2<\infty\}/\sim,
\]

где

\[
X\sim Y
\quad\Longleftrightarrow\quad
X=Y\ \text{почти наверное}.
\]

Важно: строго элементы \(L^2\) - не отдельные функции, а классы
эквивалентности по равенству почти наверное. В устном ответе обычно
говорят ``случайная величина из \(L^2\)'', понимая этот класс.

Скалярное произведение:

\[
\langle X,Y\rangle=\mathbb E[XY].
\]

Оно корректно определено, потому что по неравенству Коши-Буняковского

\[
\mathbb E|XY|
\le
(\mathbb EX^2)^{1/2}(\mathbb EY^2)^{1/2}<\infty.
\]

Норма:

\[
\|X\|_2=(\mathbb EX^2)^{1/2}.
\]

Расстояние:

\[
d(X,Y)=\|X-Y\|_2
=
\left(\mathbb E[(X-Y)^2]\right)^{1/2}.
\]

То есть близость в \(L^2\) - это малая среднеквадратичная ошибка.

Пространство \(L^2\) полно, поэтому это гильбертово пространство. В
вероятности это фундаментально: задачи прогноза, регрессии и фильтрации
можно формулировать как задачи ортогонального проектирования.

\subsubsection{2.7. Ортогональность и
ковариация}\label{ux43eux440ux442ux43eux433ux43eux43dux430ux43bux44cux43dux43eux441ux442ux44c-ux438-ux43aux43eux432ux430ux440ux438ux430ux446ux438ux44f}

Ортогональность в \(L^2\):

\[
X\perp Y
\quad\Longleftrightarrow\quad
\langle X,Y\rangle=\mathbb E[XY]=0.
\]

Если \(X\) и \(Y\) центрированы, то есть

\[
\mathbb EX=\mathbb EY=0,
\]

то

\[
\operatorname{Cov}(X,Y)=\mathbb E[XY]=\langle X,Y\rangle.
\]

В общем случае удобнее работать с центрированными величинами

\[
X^0=X-\mathbb EX,
\qquad
Y^0=Y-\mathbb EY.
\]

Тогда

\[
\operatorname{Cov}(X,Y)
=
\mathbb E[X^0Y^0]
=
\langle X^0,Y^0\rangle.
\]

Отсюда корреляция в Билете 14 имеет геометрический смысл косинуса угла:

\[
\rho(X,Y)
=
\frac{\langle X^0,Y^0\rangle}{\|X^0\|_2\|Y^0\|_2}.
\]

\subsubsection{2.8. Ортогональные
проекторы}\label{ux43eux440ux442ux43eux433ux43eux43dux430ux43bux44cux43dux44bux435-ux43fux440ux43eux435ux43aux442ux43eux440ux44b}

Теперь проекторы. Пусть \(H\) - линейное подпространство гильбертова
пространства \(\mathcal H\). Элемент \(y\in H\) называется ортогональной
проекцией \(x\) на \(H\), если

\[
x-y\perp H,
\]

то есть

\[
\langle x-y,h\rangle=0
\]

для всех \(h\in H\).

Главная теорема о проекции: если \(H\) - замкнутое линейное
подпространство гильбертова пространства \(\mathcal H\), то для любого
\(x\in\mathcal H\) существует единственный элемент \(P_Hx\in H\), такой
что

\[
x-P_Hx\perp H.
\]

Этот элемент является ближайшим к \(x\) элементом из \(H\):

\[
\|x-P_Hx\|
=
\inf_{h\in H}\|x-h\|.
\]

Более того, для любого \(h\in H\):

\[
\|x-h\|^2
=
\|x-P_Hx\|^2+\|P_Hx-h\|^2.
\]

Это просто теорема Пифагора, потому что

\[
x-h=(x-P_Hx)+(P_Hx-h),
\]

а две части ортогональны.

Отображение

\[
P_H:\mathcal H\to H
\]

называется ортогональным проектором. Оно обладает свойствами:

\begin{enumerate}
\def\labelenumi{\arabic{enumi}.}
\tightlist
\item
  линейность:
\end{enumerate}

\[
P_H(ax+by)=aP_Hx+bP_Hy;
\]

\begin{enumerate}
\def\labelenumi{\arabic{enumi}.}
\setcounter{enumi}{1}
\tightlist
\item
  идемпотентность:
\end{enumerate}

\[
P_H^2=P_H;
\]

\begin{enumerate}
\def\labelenumi{\arabic{enumi}.}
\setcounter{enumi}{2}
\tightlist
\item
  самосопряженность:
\end{enumerate}

\[
\langle P_Hx,y\rangle=\langle x,P_Hy\rangle;
\]

\begin{enumerate}
\def\labelenumi{\arabic{enumi}.}
\setcounter{enumi}{3}
\tightlist
\item
  сжатие нормы:
\end{enumerate}

\[
\|P_Hx\|\le\|x\|;
\]

\begin{enumerate}
\def\labelenumi{\arabic{enumi}.}
\setcounter{enumi}{4}
\tightlist
\item
  разложение:
\end{enumerate}

\[
x=P_Hx+(x-P_Hx),
\qquad
P_Hx\in H,
\quad
x-P_Hx\in H^\perp.
\]

Условие замкнутости \(H\) важно. Если подпространство не замкнуто,
ближайший элемент может не существовать внутри самого \(H\).

\subsubsection{2.9. Матрица Грама и линейная
регрессия}\label{ux43cux430ux442ux440ux438ux446ux430-ux433ux440ux430ux43cux430-ux438-ux43bux438ux43dux435ux439ux43dux430ux44f-ux440ux435ux433ux440ux435ux441ux441ux438ux44f}

В конечномерном случае проекция вычисляется через матрицу Грама. Пусть
\(h_1,\ldots,h_n\) - линейно независимые элементы гильбертова
пространства, а

\[
H=\operatorname{span}(h_1,\ldots,h_n).
\]

Ищем проекцию в виде

\[
P_Hx=\sum_{j=1}^n a_jh_j.
\]

Условие ортогональности ошибки:

\[
x-\sum_{j=1}^n a_jh_j\perp h_i,
\qquad i=1,\ldots,n.
\]

Значит коэффициенты удовлетворяют нормальным уравнениям:

\[
\sum_{j=1}^n a_j\langle h_j,h_i\rangle
=
\langle x,h_i\rangle,
\qquad i=1,\ldots,n.
\]

Матрица

\[
G=(\langle h_i,h_j\rangle)_{i,j=1}^n
\]

называется матрицей Грама. Если \(h_1,\ldots,h_n\) линейно независимы,
то \(G\) положительно определена и обратима. В ортонормированном случае

\[
\langle h_i,h_j\rangle=\delta_{ij},
\]

и формула упрощается:

\[
P_Hx=\sum_{j=1}^n \langle x,h_j\rangle h_j.
\]

В вероятности это дает линейную регрессию. Пусть
\(X,Y_1,\ldots,Y_n\in L^2\), и мы хотим приблизить \(X\) линейной
комбинацией \(Y_1,\ldots,Y_n\):

\[
\widehat X=\sum_{j=1}^n a_jY_j.
\]

Оптимальная в среднеквадратичном смысле оценка минимизирует

\[
\mathbb E\left(X-\sum_{j=1}^n a_jY_j\right)^2.
\]

Она является ортогональной проекцией \(X\) на

\[
H=\operatorname{span}(Y_1,\ldots,Y_n).
\]

Нормальные уравнения:

\[
\mathbb E\left[\left(X-\sum_{j=1}^n a_jY_j\right)Y_i\right]=0,
\qquad i=1,\ldots,n.
\]

То есть ошибка оптимального прогноза ортогональна всем используемым
регрессорам.

Если работают с центрированными величинами, то матрица Грама становится
ковариационной матрицей:

\[
G_{ij}
=
\langle Y_i^0,Y_j^0\rangle
=
\operatorname{Cov}(Y_i,Y_j).
\]

Поэтому ковариационные матрицы всегда неотрицательно определены:

\[
a^TGa
=
\left\|\sum_{i=1}^n a_iY_i^0\right\|_2^2
\ge0.
\]

\subsubsection{2.10. Условное ожидание как ортогональная
проекция}\label{ux443ux441ux43bux43eux432ux43dux43eux435-ux43eux436ux438ux434ux430ux43dux438ux435-ux43aux430ux43a-ux43eux440ux442ux43eux433ux43eux43dux430ux43bux44cux43dux430ux44f-ux43fux440ux43eux435ux43aux446ux438ux44f}

Теперь условное ожидание как проектор. Пусть \(X\in L^2\) и
\(\mathcal G\subset\mathcal F\) - под-\(\sigma\)-алгебра. Рассмотрим

\[
L^2(\mathcal G)
=
\{Y\in L^2(\Omega,\mathcal F,P):Y\ \mathcal G\text{-измерима}\}.
\]

Это замкнутое линейное подпространство \(L^2(\mathcal F)\). Тогда

\[
\mathbb E[X\mid\mathcal G]
\]

является ортогональной проекцией \(X\) на \(L^2(\mathcal G)\):

\[
P_{L^2(\mathcal G)}X=\mathbb E[X\mid\mathcal G].
\]

Почему? По определению условного ожидания для любого \(G\in\mathcal G\):

\[
\int_G(X-\mathbb E[X\mid\mathcal G])\,dP=0.
\]

Это означает ортогональность ошибки индикаторам \(\mathbf 1_G\). А через
линейность и предельный переход ошибка ортогональна всем ограниченным
\(\mathcal G\)-измеримым величинам, а затем всем величинам из
\(L^2(\mathcal G)\). Значит

\[
X-\mathbb E[X\mid\mathcal G]\perp L^2(\mathcal G).
\]

Отсюда следует оптимальное свойство условного ожидания:

\[
\mathbb E[(X-\mathbb E[X\mid\mathcal G])^2]
=
\min_{Y\in L^2(\mathcal G)}
\mathbb E[(X-Y)^2].
\]

То есть условное ожидание - лучший среднеквадратичный прогноз \(X\) по
информации \(\mathcal G\).

\subsubsection{2.11.
Заключение}\label{ux437ux430ux43aux43bux44eux447ux435ux43dux438ux435}

Итак, билет строит мост: порядок и алгебраические операции задают общий
язык; \(L^2\) превращает случайные величины в векторы; ковариации
становятся скалярными произведениями; условные ожидания и линейные
оценки становятся ортогональными проекциями.

\subsection{3. Основные
определения}\label{ux43eux441ux43dux43eux432ux43dux44bux435-ux43eux43fux440ux435ux434ux435ux43bux435ux43dux438ux44f}

\subsubsection{Частично упорядоченное
множество}\label{ux447ux430ux441ux442ux438ux447ux43dux43e-ux443ux43fux43eux440ux44fux434ux43eux447ux435ux43dux43dux43eux435-ux43cux43dux43eux436ux435ux441ux442ux432ux43e}

Множество \(M\) с отношением \(\le\), которое рефлексивно,
антисимметрично и транзитивно.

\subsubsection{Линейный
порядок}\label{ux43bux438ux43dux435ux439ux43dux44bux439-ux43fux43eux440ux44fux434ux43eux43a}

Частичный порядок, в котором любые два элемента сравнимы.

\subsubsection{Верхняя и нижняя
грань}\label{ux432ux435ux440ux445ux43dux44fux44f-ux438-ux43dux438ux436ux43dux44fux44f-ux433ux440ux430ux43dux44c}

Для \(A\subset M\) элемент \(u\) - верхняя грань, если \(a\le u\) для
всех \(a\in A\). Нижняя грань определяется аналогично.

\subsubsection{Супремум и
инфимум}\label{ux441ux443ux43fux440ux435ux43cux443ux43c-ux438-ux438ux43dux444ux438ux43cux443ux43c}

Наименьшая верхняя грань называется супремумом, наибольшая нижняя грань
- инфимумом.

\subsubsection{Алгебраическая
структура}\label{ux430ux43bux433ux435ux431ux440ux430ux438ux447ux435ux441ux43aux430ux44f-ux441ux442ux440ux443ux43aux442ux443ux440ux430}

Множество с заданными операциями, удовлетворяющими аксиомам. Примеры:
группа, кольцо, поле, линейное пространство.

\subsubsection{Линейное
пространство}\label{ux43bux438ux43dux435ux439ux43dux43eux435-ux43fux440ux43eux441ux442ux440ux430ux43dux441ux442ux432ux43e}

Множество \(V\), где определены сложение векторов и умножение на
скаляры, причем выполнены стандартные аксиомы линейности.

\subsubsection{Линейное
подпространство}\label{ux43bux438ux43dux435ux439ux43dux43eux435-ux43fux43eux434ux43fux440ux43eux441ux442ux440ux430ux43dux441ux442ux432ux43e}

Подмножество \(H\subset V\), замкнутое относительно линейных комбинаций:

\[
x,y\in H,\ a,b\in\mathbb R
\quad\Rightarrow\quad
ax+by\in H.
\]

\subsubsection{Скалярное
произведение}\label{ux441ux43aux430ux43bux44fux440ux43dux43eux435-ux43fux440ux43eux438ux437ux432ux435ux434ux435ux43dux438ux435}

Функция \(\langle x,y\rangle\), линейная по одному аргументу,
симметричная в вещественном случае и положительно определенная.

\subsubsection{Норма, порожденная скалярным
произведением}\label{ux43dux43eux440ux43cux430-ux43fux43eux440ux43eux436ux434ux435ux43dux43dux430ux44f-ux441ux43aux430ux43bux44fux440ux43dux44bux43c-ux43fux440ux43eux438ux437ux432ux435ux434ux435ux43dux438ux435ux43c}

\[
\|x\|=\sqrt{\langle x,x\rangle}.
\]

\subsubsection{Гильбертово
пространство}\label{ux433ux438ux43bux44cux431ux435ux440ux442ux43eux432ux43e-ux43fux440ux43eux441ux442ux440ux430ux43dux441ux442ux432ux43e}

Полное линейное пространство со скалярным произведением.

\subsubsection{\texorpdfstring{Пространство
\(L^2\)}{Пространство L\^{}2}}\label{ux43fux440ux43eux441ux442ux440ux430ux43dux441ux442ux432ux43e-l2-1}

\[
L^2(\Omega,\mathcal F,P)=\{X:\mathbb EX^2<\infty\}/\sim,
\]

где \(X\sim Y\), если \(X=Y\) почти наверное.

\subsubsection{\texorpdfstring{Скалярное произведение в
\(L^2\)}{Скалярное произведение в L\^{}2}}\label{ux441ux43aux430ux43bux44fux440ux43dux43eux435-ux43fux440ux43eux438ux437ux432ux435ux434ux435ux43dux438ux435-ux432-l2}

\[
\langle X,Y\rangle=\mathbb E[XY].
\]

Для комплексных величин обычно берут

\[
\langle X,Y\rangle=\mathbb E[X\overline Y]
\]

или сопряжение в другом аргументе, в зависимости от конвенции.

\subsubsection{Ортогональность}\label{ux43eux440ux442ux43eux433ux43eux43dux430ux43bux44cux43dux43eux441ux442ux44c}

\[
x\perp y
\quad\Longleftrightarrow\quad
\langle x,y\rangle=0.
\]

В \(L^2\):

\[
X\perp Y
\quad\Longleftrightarrow\quad
\mathbb E[XY]=0.
\]

\subsubsection{Ортогональное
дополнение}\label{ux43eux440ux442ux43eux433ux43eux43dux430ux43bux44cux43dux43eux435-ux434ux43eux43fux43eux43bux43dux435ux43dux438ux435}

\[
H^\perp=\{x:\langle x,h\rangle=0\ \text{для всех }h\in H\}.
\]

\subsubsection{Ортогональная
проекция}\label{ux43eux440ux442ux43eux433ux43eux43dux430ux43bux44cux43dux430ux44f-ux43fux440ux43eux435ux43aux446ux438ux44f}

Элемент \(P_Hx\in H\), такой что

\[
x-P_Hx\perp H.
\]

\subsubsection{Ортогональный
проектор}\label{ux43eux440ux442ux43eux433ux43eux43dux430ux43bux44cux43dux44bux439-ux43fux440ux43eux435ux43aux442ux43eux440}

Отображение \(P_H:x\mapsto P_Hx\) на замкнутое подпространство \(H\).

\subsubsection{Матрица
Грама}\label{ux43cux430ux442ux440ux438ux446ux430-ux433ux440ux430ux43cux430}

Для \(h_1,\ldots,h_n\):

\[
G=(\langle h_i,h_j\rangle)_{i,j=1}^n.
\]

\subsubsection{Ортонормированная
система}\label{ux43eux440ux442ux43eux43dux43eux440ux43cux438ux440ux43eux432ux430ux43dux43dux430ux44f-ux441ux438ux441ux442ux435ux43cux430}

Система \((e_i)\), для которой

\[
\langle e_i,e_j\rangle=\delta_{ij}.
\]

\subsection{4. Основные теоремы и
свойства}\label{ux43eux441ux43dux43eux432ux43dux44bux435-ux442ux435ux43eux440ux435ux43cux44b-ux438-ux441ux432ux43eux439ux441ux442ux432ux430}

\subsubsection{\texorpdfstring{Теорема 1. \(L^2\) - гильбертово
пространство}{Теорема 1. L\^{}2 - гильбертово пространство}}\label{ux442ux435ux43eux440ux435ux43cux430-1.-l2---ux433ux438ux43bux44cux431ux435ux440ux442ux43eux432ux43e-ux43fux440ux43eux441ux442ux440ux430ux43dux441ux442ux432ux43e}

Пространство случайных величин с конечным вторым моментом и скалярным
произведением

\[
\langle X,Y\rangle=\mathbb E[XY]
\]

является гильбертовым пространством после отождествления случайных
величин, равных почти наверное.

\subsubsection{Условие
корректности}\label{ux443ux441ux43bux43eux432ux438ux435-ux43aux43eux440ux440ux435ux43aux442ux43dux43eux441ux442ux438}

Если \(X,Y\in L^2\), то \(\mathbb E|XY|<\infty\) по Коши-Буняковскому.

\subsubsection{Теорема 2.
Коши-Буняковский}\label{ux442ux435ux43eux440ux435ux43cux430-2.-ux43aux43eux448ux438-ux431ux443ux43dux44fux43aux43eux432ux441ux43aux438ux439}

В любом пространстве со скалярным произведением:

\[
|\langle x,y\rangle|
\le
\|x\|\|y\|.
\]

В \(L^2\):

\[
|\mathbb E[XY]|
\le
(\mathbb EX^2)^{1/2}(\mathbb EY^2)^{1/2}.
\]

\subsubsection{Теорема 3. Проекция на замкнутое
подпространство}\label{ux442ux435ux43eux440ux435ux43cux430-3.-ux43fux440ux43eux435ux43aux446ux438ux44f-ux43dux430-ux437ux430ux43cux43aux43dux443ux442ux43eux435-ux43fux43eux434ux43fux440ux43eux441ux442ux440ux430ux43dux441ux442ux432ux43e}

Если \(H\) - замкнутое линейное подпространство гильбертова пространства
\(\mathcal H\), то для любого \(x\in\mathcal H\) существует единственный
\(P_Hx\in H\), такой что

\[
x-P_Hx\perp H.
\]

При этом

\[
\|x-P_Hx\|
=
\min_{h\in H}\|x-h\|.
\]

\subsubsection{Свойство 4. Пифагорова формула для
проекции}\label{ux441ux432ux43eux439ux441ux442ux432ux43e-4.-ux43fux438ux444ux430ux433ux43eux440ux43eux432ux430-ux444ux43eux440ux43cux443ux43bux430-ux434ux43bux44f-ux43fux440ux43eux435ux43aux446ux438ux438}

Для любого \(h\in H\):

\[
\|x-h\|^2
=
\|x-P_Hx\|^2+\|P_Hx-h\|^2.
\]

\subsubsection{Свойство 5. Свойства ортогонального
проектора}\label{ux441ux432ux43eux439ux441ux442ux432ux43e-5.-ux441ux432ux43eux439ux441ux442ux432ux430-ux43eux440ux442ux43eux433ux43eux43dux430ux43bux44cux43dux43eux433ux43e-ux43fux440ux43eux435ux43aux442ux43eux440ux430}

Ортогональный проектор \(P_H\):

\[
P_H^2=P_H,
\]

линеен,

\[
\langle P_Hx,y\rangle=\langle x,P_Hy\rangle,
\]

и

\[
\|P_Hx\|\le\|x\|.
\]

\subsubsection{Теорема 6. Нормальные уравнения конечномерной
проекции}\label{ux442ux435ux43eux440ux435ux43cux430-6.-ux43dux43eux440ux43cux430ux43bux44cux43dux44bux435-ux443ux440ux430ux432ux43dux435ux43dux438ux44f-ux43aux43eux43dux435ux447ux43dux43eux43cux435ux440ux43dux43eux439-ux43fux440ux43eux435ux43aux446ux438ux438}

Если \(H=\operatorname{span}(h_1,\ldots,h_n)\) и

\[
P_Hx=\sum_{j=1}^n a_jh_j,
\]

то коэффициенты \(a_j\) определяются системой

\[
\sum_{j=1}^n a_j\langle h_j,h_i\rangle
=
\langle x,h_i\rangle,
\qquad i=1,\ldots,n.
\]

\subsubsection{Свойство 7. Проекция на ортонормированную
систему}\label{ux441ux432ux43eux439ux441ux442ux432ux43e-7.-ux43fux440ux43eux435ux43aux446ux438ux44f-ux43dux430-ux43eux440ux442ux43eux43dux43eux440ux43cux438ux440ux43eux432ux430ux43dux43dux443ux44e-ux441ux438ux441ux442ux435ux43cux443}

Если \(e_1,\ldots,e_n\) ортонормированы, то

\[
P_Hx=\sum_{j=1}^n\langle x,e_j\rangle e_j.
\]

\subsubsection{Свойство 8. Матрица Грама неотрицательно
определена}\label{ux441ux432ux43eux439ux441ux442ux432ux43e-8.-ux43cux430ux442ux440ux438ux446ux430-ux433ux440ux430ux43cux430-ux43dux435ux43eux442ux440ux438ux446ux430ux442ux435ux43bux44cux43dux43e-ux43eux43fux440ux435ux434ux435ux43bux435ux43dux430}

Для любых \(a_1,\ldots,a_n\):

\[
\sum_{i,j=1}^n a_i a_j\langle h_i,h_j\rangle
=
\left\|\sum_{i=1}^n a_ih_i\right\|^2
\ge0.
\]

Если \(h_1,\ldots,h_n\) линейно независимы, то матрица Грама
положительно определена.

\subsubsection{Свойство 9. Ковариационная матрица как матрица
Грама}\label{ux441ux432ux43eux439ux441ux442ux432ux43e-9.-ux43aux43eux432ux430ux440ux438ux430ux446ux438ux43eux43dux43dux430ux44f-ux43cux430ux442ux440ux438ux446ux430-ux43aux430ux43a-ux43cux430ux442ux440ux438ux446ux430-ux433ux440ux430ux43cux430}

Для центрированных величин \(X_1,\ldots,X_n\in L^2\):

\[
G_{ij}=\mathbb E[X_iX_j]
\]

\begin{itemize}
\tightlist
\item
  матрица Грама. Поэтому она неотрицательно определена.
\end{itemize}

Для нецентрированных величин ковариационная матрица является матрицей
Грама центрированных величин:

\[
\Sigma_{ij}
=
\operatorname{Cov}(X_i,X_j)
=
\langle X_i-\mathbb EX_i,\ X_j-\mathbb EX_j\rangle.
\]

\subsubsection{Теорема 10. Условное ожидание как ортогональная
проекция}\label{ux442ux435ux43eux440ux435ux43cux430-10.-ux443ux441ux43bux43eux432ux43dux43eux435-ux43eux436ux438ux434ux430ux43dux438ux435-ux43aux430ux43a-ux43eux440ux442ux43eux433ux43eux43dux430ux43bux44cux43dux430ux44f-ux43fux440ux43eux435ux43aux446ux438ux44f}

Если \(X\in L^2\) и \(\mathcal G\subset\mathcal F\), то

\[
\mathbb E[X\mid\mathcal G]
=
P_{L^2(\mathcal G)}X.
\]

Эквивалентно:

\[
X-\mathbb E[X\mid\mathcal G]\perp L^2(\mathcal G).
\]

И оптимальное свойство:

\[
\mathbb E[(X-\mathbb E[X\mid\mathcal G])^2]
=
\min_{Y\in L^2(\mathcal G)}
\mathbb E[(X-Y)^2].
\]

\subsection{5. Примеры}\label{ux43fux440ux438ux43cux435ux440ux44b}

\subsubsection{Пример 1. Евклидово
пространство}\label{ux43fux440ux438ux43cux435ux440-1.-ux435ux432ux43aux43bux438ux434ux43eux432ux43e-ux43fux440ux43eux441ux442ux440ux430ux43dux441ux442ux432ux43e}

В \(\mathbb R^n\):

\[
\langle x,y\rangle=x^Ty,
\qquad
\|x\|=\sqrt{x_1^2+\cdots+x_n^2}.
\]

Проекция точки на прямую или плоскость - обычная ортогональная проекция
из геометрии.

\subsubsection{\texorpdfstring{Пример 2. Пространство
\(L^2\)}{Пример 2. Пространство L\^{}2}}\label{ux43fux440ux438ux43cux435ux440-2.-ux43fux440ux43eux441ux442ux440ux430ux43dux441ux442ux432ux43e-l2}

Пусть \(X,Y\in L^2\). Тогда

\[
\langle X,Y\rangle=\mathbb E[XY].
\]

Если \(\mathbb EX=\mathbb EY=0\), то

\[
X\perp Y
\quad\Longleftrightarrow\quad
\operatorname{Cov}(X,Y)=0.
\]

То есть ортогональность центрированных величин - это
некоррелированность.

\subsubsection{Пример 3. Проекция на
константы}\label{ux43fux440ux438ux43cux435ux440-3.-ux43fux440ux43eux435ux43aux446ux438ux44f-ux43dux430-ux43aux43eux43dux441ux442ux430ux43dux442ux44b}

Пусть \(H\) - подпространство констант:

\[
H=\{c:\ c\in\mathbb R\}.
\]

Проекция \(X\in L^2\) на \(H\) равна

\[
P_HX=\mathbb EX.
\]

Действительно,

\[
\mathbb E[X-\mathbb EX]=0,
\]

то есть ошибка ортогональна всем константам.

\subsubsection{Пример 4. Условное ожидание по
разбиению}\label{ux43fux440ux438ux43cux435ux440-4.-ux443ux441ux43bux43eux432ux43dux43eux435-ux43eux436ux438ux434ux430ux43dux438ux435-ux43fux43e-ux440ux430ux437ux431ux438ux435ux43dux438ux44e}

Пусть \(\mathcal G=\sigma(D_1,\ldots,D_n)\), где \(D_i\) - атомы
разбиения и \(P(D_i)>0\). Тогда

\[
\mathbb E[X\mid\mathcal G]
=
\sum_{i=1}^n
\frac{\int_{D_i}X\,dP}{P(D_i)}\mathbf 1_{D_i}.
\]

Это проекция \(X\) на подпространство функций, постоянных на каждом
\(D_i\).

\subsubsection{Пример 5. Линейная регрессия на одну
величину}\label{ux43fux440ux438ux43cux435ux440-5.-ux43bux438ux43dux435ux439ux43dux430ux44f-ux440ux435ux433ux440ux435ux441ux441ux438ux44f-ux43dux430-ux43eux434ux43dux443-ux432ux435ux43bux438ux447ux438ux43dux443}

Пусть \(X,Y\in L^2\), \(\mathbb EY=0\), \(\mathbb EY^2>0\). Ищем лучшую
оценку \(X\) вида \(aY\). Минимизируем

\[
\mathbb E(X-aY)^2.
\]

Условие ортогональности:

\[
\mathbb E[(X-aY)Y]=0.
\]

Отсюда

\[
a=\frac{\mathbb E[XY]}{\mathbb E[Y^2]}.
\]

Если \(X,Y\) центрированы, то

\[
a=\frac{\operatorname{Cov}(X,Y)}{\operatorname{Var}Y}.
\]

\subsubsection{Пример 6. Матрица Грама и ковариационная
матрица}\label{ux43fux440ux438ux43cux435ux440-6.-ux43cux430ux442ux440ux438ux446ux430-ux433ux440ux430ux43cux430-ux438-ux43aux43eux432ux430ux440ux438ux430ux446ux438ux43eux43dux43dux430ux44f-ux43cux430ux442ux440ux438ux446ux430}

Пусть \(X=(X_1,\ldots,X_n)\), где все \(X_i\) центрированы и имеют
вторые моменты. Тогда

\[
G_{ij}=\mathbb E[X_iX_j]
\]

является матрицей Грама. Для любого \(a\in\mathbb R^n\):

\[
a^TGa
=
\mathbb E\left(\sum_{i=1}^n a_iX_i\right)^2
\ge0.
\]

Поэтому ковариационная матрица неотрицательно определена.

\subsubsection{Пример 7. Ортогональность не равна
независимости}\label{ux43fux440ux438ux43cux435ux440-7.-ux43eux440ux442ux43eux433ux43eux43dux430ux43bux44cux43dux43eux441ux442ux44c-ux43dux435-ux440ux430ux432ux43dux430-ux43dux435ux437ux430ux432ux438ux441ux438ux43cux43eux441ux442ux438}

Пусть \(X\) равномерна на \([-1,1]\), а

\[
Y=X^2.
\]

Тогда \(Y\) полностью определяется по \(X\), значит независимости нет.
Но

\[
\mathbb E[X]=0,
\qquad
\mathbb E[XY]=\mathbb E[X^3]=0,
\]

поэтому центрированные величины могут быть ортогональны или
некоррелированы без независимости.

\subsection{6. Схемы доказательств /
обоснований}\label{ux441ux445ux435ux43cux44b-ux434ux43eux43aux430ux437ux430ux442ux435ux43bux44cux441ux442ux432-ux43eux431ux43eux441ux43dux43eux432ux430ux43dux438ux439}

\textbf{Почему \(L^2\) - пространство со скалярным произведением.}

Для \(X,Y\in L^2\) произведение \(XY\) интегрируемо по
Коши-Буняковскому:

\[
\mathbb E|XY|
\le
(\mathbb EX^2)^{1/2}(\mathbb EY^2)^{1/2}.
\]

Линейность и симметрия следуют из линейности математического ожидания.
Положительность:

\[
\langle X,X\rangle=\mathbb EX^2\ge0.
\]

Если \(\mathbb EX^2=0\), то \(X=0\) почти наверное, поэтому в \(L^2\)
это нулевой элемент.

\textbf{Идея полноты \(L^2\).}

Если \((X_n)\) фундаментальна по норме \(\|\cdot\|_2\), то существует
\(X\in L^2\), к которому \(X_n\) сходится в среднем квадратичном:

\[
\mathbb E[(X_n-X)^2]\to0.
\]

Полное доказательство относится к теории \(L^p\), но для экзамена важно
знать результат: \(L^2\) полно, значит это гильбертово пространство.

\textbf{Доказательство теоремы о проекции: идея.}

Нужно минимизировать функцию

\[
h\mapsto\|x-h\|^2
\]

по \(h\in H\). В гильбертовом пространстве замкнутость \(H\)
обеспечивает существование ближайшего элемента. Если \(y=P_Hx\) -
минимум, то для любого направления \(u\in H\) функция

\[
\varphi(t)=\|x-y-tu\|^2
\]

имеет минимум при \(t=0\). Поэтому

\[
\varphi'(0)=-2\langle x-y,u\rangle=0.
\]

Значит

\[
x-y\perp H.
\]

Обратно, если \(x-y\perp H\), то для любого \(h\in H\)

\[
x-h=(x-y)+(y-h)
\]

и по Пифагору

\[
\|x-h\|^2=\|x-y\|^2+\|y-h\|^2\ge\|x-y\|^2.
\]

Значит \(y\) - ближайший элемент.

\textbf{Доказательство нормальных уравнений.}

Пусть

\[
\widehat x=\sum_{j=1}^n a_jh_j.
\]

Условие проекции:

\[
x-\widehat x\perp h_i
\]

для каждого \(i\). Поэтому

\[
\left\langle x-\sum_{j=1}^n a_jh_j,\ h_i\right\rangle=0.
\]

Раскрывая линейность, получаем

\[
\sum_{j=1}^n a_j\langle h_j,h_i\rangle=\langle x,h_i\rangle.
\]

\textbf{Почему условное ожидание - проекция.}

Пусть \(Y=\mathbb E[X\mid\mathcal G]\). По определению условного
ожидания

\[
\int_G(X-Y)\,dP=0
\]

для всех \(G\in\mathcal G\). Это значит

\[
\mathbb E[(X-Y)\mathbf 1_G]=0.
\]

Линейные комбинации индикаторов дают простые \(\mathcal G\)-измеримые
функции. Предельным переходом получаем

\[
\mathbb E[(X-Y)Z]=0
\]

для всех \(Z\in L^2(\mathcal G)\). Значит

\[
X-Y\perp L^2(\mathcal G),
\]

и \(Y\) является ортогональной проекцией.

\subsection{7. Что написать на
доске}\label{ux447ux442ux43e-ux43dux430ux43fux438ux441ux430ux442ux44c-ux43dux430-ux434ux43eux441ux43aux435}

\begin{enumerate}
\def\labelenumi{\arabic{enumi}.}
\tightlist
\item
  Частичный порядок:
\end{enumerate}

\[
a\le a,\quad
a\le b,\ b\le a\Rightarrow a=b,\quad
a\le b,\ b\le c\Rightarrow a\le c.
\]

\begin{enumerate}
\def\labelenumi{\arabic{enumi}.}
\setcounter{enumi}{1}
\tightlist
\item
  \(L^2\):
\end{enumerate}

\[
L^2=\{X:\mathbb EX^2<\infty\}/\sim.
\]

\begin{enumerate}
\def\labelenumi{\arabic{enumi}.}
\setcounter{enumi}{2}
\tightlist
\item
  Скалярное произведение:
\end{enumerate}

\[
\langle X,Y\rangle=\mathbb E[XY],
\qquad
\|X\|_2=(\mathbb EX^2)^{1/2}.
\]

\begin{enumerate}
\def\labelenumi{\arabic{enumi}.}
\setcounter{enumi}{3}
\tightlist
\item
  Центрированные величины:
\end{enumerate}

\[
\operatorname{Cov}(X,Y)
=
\langle X-\mathbb EX,\ Y-\mathbb EY\rangle.
\]

\begin{enumerate}
\def\labelenumi{\arabic{enumi}.}
\setcounter{enumi}{4}
\tightlist
\item
  Проектор:
\end{enumerate}

\[
P_Hx\in H,
\qquad
x-P_Hx\perp H.
\]

\begin{enumerate}
\def\labelenumi{\arabic{enumi}.}
\setcounter{enumi}{5}
\tightlist
\item
  Нормальные уравнения:
\end{enumerate}

\[
\sum_j a_j\langle h_j,h_i\rangle=\langle x,h_i\rangle.
\]

\begin{enumerate}
\def\labelenumi{\arabic{enumi}.}
\setcounter{enumi}{6}
\tightlist
\item
  Условное ожидание как проекция:
\end{enumerate}

\[
\mathbb E[X\mid\mathcal G]
=
P_{L^2(\mathcal G)}X.
\]

\begin{enumerate}
\def\labelenumi{\arabic{enumi}.}
\setcounter{enumi}{7}
\tightlist
\item
  Матрица Грама:
\end{enumerate}

\[
G_{ij}=\langle h_i,h_j\rangle,
\qquad
a^TGa=\left\|\sum_i a_ih_i\right\|^2\ge0.
\]

\subsection{8. Типичные дополнительные
вопросы}\label{ux442ux438ux43fux438ux447ux43dux44bux435-ux434ux43eux43fux43eux43bux43dux438ux442ux435ux43bux44cux43dux44bux435-ux432ux43eux43fux440ux43eux441ux44b}

\textbf{Вопрос. Что такое частичный порядок?}

Это отношение \(\le\), которое рефлексивно, антисимметрично и
транзитивно. Пример: включение множеств \(A\subset B\).

\textbf{Вопрос. Чем частичный порядок отличается от линейного?}

В линейном порядке любые два элемента сравнимы. В частичном - не
обязательно. Например, два события могут не быть вложены друг в друга.

\textbf{Вопрос. Какие алгебраические структуры важны в вероятности?}

Линейные пространства случайных величин, пространства \(L^p\), алгебры и
\(\sigma\)-алгебры событий, матрицы с операциями сложения и умножения.

\textbf{Вопрос. Почему \(L^2\) - линейное пространство?}

Если \(X,Y\in L^2\), то \(aX+bY\in L^2\), потому что квадрат линейной
комбинации оценивается через \(X^2\) и \(Y^2\).

\textbf{Вопрос. Почему \(\mathbb E[XY]\) корректно определено для
\(X,Y\in L^2\)?}

По неравенству Коши-Буняковского:

\[
\mathbb E|XY|\le(\mathbb EX^2)^{1/2}(\mathbb EY^2)^{1/2}.
\]

\textbf{Вопрос. Что значит полнота \(L^2\)?}

Всякая фундаментальная по среднеквадратичной норме последовательность
случайных величин сходится в \(L^2\) к случайной величине из \(L^2\).

\textbf{Вопрос. Что такое ортогональность в \(L^2\)?}

\[
X\perp Y
\quad\Longleftrightarrow\quad
\mathbb E[XY]=0.
\]

Для центрированных величин это то же самое, что некоррелированность.

\textbf{Вопрос. Ортогональность означает независимость?}

Нет. Ортогональность означает нулевое скалярное произведение, то есть
для центрированных величин нулевую ковариацию. Независимость сильнее.

\textbf{Вопрос. Зачем нужна замкнутость подпространства в теореме о
проекции?}

Чтобы ближайший элемент действительно лежал в \(H\). Без замкнутости
минимум может достигаться только в замыкании \(H\).

\textbf{Вопрос. Как найти проекцию на конечномерное подпространство?}

Записать

\[
P_Hx=\sum_j a_jh_j
\]

и решить нормальные уравнения:

\[
\sum_j a_j\langle h_j,h_i\rangle=\langle x,h_i\rangle.
\]

\textbf{Вопрос. Что такое матрица Грама?}

\[
G_{ij}=\langle h_i,h_j\rangle.
\]

Она неотрицательно определена, потому что

\[
a^TGa=\left\|\sum_i a_ih_i\right\|^2\ge0.
\]

\textbf{Вопрос. Почему ковариационная матрица неотрицательно
определена?}

Она является матрицей Грама центрированных случайных величин:

\[
\Sigma_{ij}
=
\langle X_i-\mathbb EX_i,\ X_j-\mathbb EX_j\rangle.
\]

\textbf{Вопрос. Почему условное ожидание - проекция?}

Потому что ошибка

\[
X-\mathbb E[X\mid\mathcal G]
\]

ортогональна всем \(\mathcal G\)-измеримым величинам из \(L^2\).

\textbf{Вопрос. Какой оптимальный смысл условного ожидания?}

Это лучший среднеквадратичный прогноз:

\[
\mathbb E[(X-\mathbb E[X\mid\mathcal G])^2]
=
\min_{Y\in L^2(\mathcal G)}
\mathbb E[(X-Y)^2].
\]

\subsection{9. Типичные
ошибки}\label{ux442ux438ux43fux438ux447ux43dux44bux435-ux43eux448ux438ux431ux43aux438}

\begin{enumerate}
\def\labelenumi{\arabic{enumi}.}
\tightlist
\item
  Путать ортогональность и независимость.
\item
  Забывать, что в \(L^2\) элементы считаются с точностью до равенства
  почти наверное.
\item
  Использовать скалярное произведение \(\mathbb E[XY]\) без проверки
  вторых моментов.
\item
  Называть любое линейное подпространство подходящим для проекции,
  забывая замкнутость.
\item
  Путать матрицу вторых моментов \(\mathbb E[XX^T]\) и ковариационную
  матрицу \(\mathbb E[(X-m)(X-m)^T]\).
\item
  Забывать центрирование при связи ковариации со скалярным
  произведением.
\item
  Считать, что условное ожидание как проекция верно только для конечных
  разбиений. На самом деле оно верно для любой под-\(\sigma\)-алгебры
  при \(X\in L^2\).
\item
  В нормальных уравнениях ставить коэффициенты без матрицы Грама, как
  будто система всегда ортонормирована.
\item
  Путать проектор \(P_H\) и произвольный линейный оператор с \(P^2=P\):
  ортогональный проектор еще самосопряжен.
\item
  Забывать, что \(L^2\)-норма измеряет среднеквадратичную ошибку, а не
  поточечное расстояние.
\end{enumerate}

\subsection{10. Связи с другими
билетами}\label{ux441ux432ux44fux437ux438-ux441-ux434ux440ux443ux433ux438ux43cux438-ux431ux438ux43bux435ux442ux430ux43cux438}

\textbf{Билет 10.} Условное ожидание определяется через интегральное
свойство, а в \(L^2\) получает геометрический смысл ортогональной
проекции.

\textbf{Билет 11.} Вторые моменты, дисперсии и ковариации вычисляются
через интегралы по распределению; здесь они становятся нормами и
скалярными произведениями.

\textbf{Билет 13.} Матрица вторых моментов и ковариационная матрица
связаны с производными характеристической функции.

\textbf{Билет 14.} Коэффициент корреляции - косинус угла между
центрированными величинами в \(L^2\).

\textbf{Билет 18.} Гауссовские векторы задаются средним и ковариационной
матрицей; ортогонализация и проекции особенно удобны в гауссовском
случае.

\textbf{Билет 24.} Корреляционная функция и эргодичность по среднему
используют \(L^2\)-оценки и дисперсии средних.

\textbf{Билет 25.} Спектральная теория стационарных последовательностей
строится в гильбертовом пространстве второго порядка.

\subsection{11. Минимум на
удовлетворительно}\label{ux43cux438ux43dux438ux43cux443ux43c-ux43dux430-ux443ux434ux43eux432ux43bux435ux442ux432ux43eux440ux438ux442ux435ux43bux44cux43dux43e}

Нужно уметь сказать:

\begin{enumerate}
\def\labelenumi{\arabic{enumi}.}
\tightlist
\item
  Частичный порядок - рефлексивность, антисимметричность,
  транзитивность.
\item
  \(L^2\) - пространство случайных величин с конечным вторым моментом:
\end{enumerate}

\[
L^2=\{X:\mathbb EX^2<\infty\}.
\]

\begin{enumerate}
\def\labelenumi{\arabic{enumi}.}
\setcounter{enumi}{2}
\tightlist
\item
  Скалярное произведение:
\end{enumerate}

\[
\langle X,Y\rangle=\mathbb E[XY].
\]

\begin{enumerate}
\def\labelenumi{\arabic{enumi}.}
\setcounter{enumi}{3}
\tightlist
\item
  Для центрированных величин:
\end{enumerate}

\[
\operatorname{Cov}(X,Y)=\langle X,Y\rangle.
\]

\begin{enumerate}
\def\labelenumi{\arabic{enumi}.}
\setcounter{enumi}{4}
\tightlist
\item
  Проекция на замкнутое подпространство:
\end{enumerate}

\[
x-P_Hx\perp H.
\]

\begin{enumerate}
\def\labelenumi{\arabic{enumi}.}
\setcounter{enumi}{5}
\tightlist
\item
  Условное ожидание в \(L^2\):
\end{enumerate}

\[
\mathbb E[X\mid\mathcal G]
=
P_{L^2(\mathcal G)}X.
\]

\subsection{12. Минимум на
хорошо}\label{ux43cux438ux43dux438ux43cux443ux43c-ux43dux430-ux445ux43eux440ux43eux448ux43e}

Помимо минимума, нужно:

\begin{enumerate}
\def\labelenumi{\arabic{enumi}.}
\tightlist
\item
  объяснить, что \(L^2\) является гильбертовым пространством;
\item
  записать норму
\end{enumerate}

\[
\|X\|_2=(\mathbb EX^2)^{1/2};
\]

\begin{enumerate}
\def\labelenumi{\arabic{enumi}.}
\setcounter{enumi}{2}
\tightlist
\item
  сформулировать теорему о проекции на замкнутое подпространство;
\item
  записать нормальные уравнения для конечномерной проекции;
\item
  объяснить матрицу Грама и ее неотрицательную определенность;
\item
  привести пример проекции на константы:
\end{enumerate}

\[
P_HX=\mathbb EX.
\]

\subsection{13. Что добавить для
отлично}\label{ux447ux442ux43e-ux434ux43eux431ux430ux432ux438ux442ux44c-ux434ux43bux44f-ux43eux442ux43bux438ux447ux43dux43e}

Для сильного ответа стоит добавить:

\begin{enumerate}
\def\labelenumi{\arabic{enumi}.}
\tightlist
\item
  различие между частичным и линейным порядком;
\item
  примеры алгебраических структур в вероятности;
\item
  доказательную идею проекторной теоремы через минимизацию
  \(\|x-h\|^2\);
\item
  Пифагорову формулу для проекции;
\item
  свойства проектора: линейность, идемпотентность, самосопряженность;
\item
  условное ожидание как лучший среднеквадратичный прогноз;
\item
  предупреждение, что ортогональность не равна независимости.
\end{enumerate}

\subsection{14. Устная версия ответа на 5
минут}\label{ux443ux441ux442ux43dux430ux44f-ux432ux435ux440ux441ux438ux44f-ux43eux442ux432ux435ux442ux430-ux43dux430-5-ux43cux438ux43dux443ux442}

Сначала я называю общий язык. Упорядоченная структура - это множество с
отношением порядка. Частичный порядок рефлексивен, антисимметричен и
транзитивен. Примеры: включение событий \(A\subset B\), порядок
случайных величин \(X\le Y\) почти наверное, порядок на неотрицательно
определенных матрицах.

Алгебраическая структура - это множество с операциями. В этом билете
главный пример - линейное пространство случайных величин. Если
\(X,Y\in L^2\), то \(aX+bY\in L^2\).

Дальше перехожу к гильбертову пространству. Рассматриваем

\[
L^2(\Omega,\mathcal F,P)=\{X:\mathbb EX^2<\infty\}/\sim.
\]

Скалярное произведение:

\[
\langle X,Y\rangle=\mathbb E[XY],
\]

норма:

\[
\|X\|_2=(\mathbb EX^2)^{1/2}.
\]

По Коши-Буняковскому \(\mathbb E[XY]\) корректно определено.
Пространство \(L^2\) полно, поэтому это гильбертово пространство.

Для центрированных величин

\[
X^0=X-\mathbb EX,\qquad Y^0=Y-\mathbb EY
\]

получаем

\[
\operatorname{Cov}(X,Y)=\langle X^0,Y^0\rangle.
\]

Значит ковариация - это скалярное произведение центрированных случайных
величин.

Потом формулирую теорему о проекции. Если \(H\) - замкнутое
подпространство гильбертова пространства, то для любого \(x\) существует
единственный \(P_Hx\in H\), такой что

\[
x-P_Hx\perp H.
\]

Это ближайший элемент из \(H\):

\[
\|x-P_Hx\|=\min_{h\in H}\|x-h\|.
\]

Проектор линеен, идемпотентен и самосопряжен.

В конечномерном случае, если

\[
P_Hx=\sum_j a_jh_j,
\]

то коэффициенты находятся из нормальных уравнений:

\[
\sum_j a_j\langle h_j,h_i\rangle=\langle x,h_i\rangle.
\]

Матрица \(\langle h_i,h_j\rangle\) - матрица Грама. Она неотрицательно
определена.

Завершаю вероятностным смыслом. Если \(X\in L^2\) и
\(\mathcal G\subset\mathcal F\), то

\[
\mathbb E[X\mid\mathcal G]
=
P_{L^2(\mathcal G)}X.
\]

То есть условное ожидание - лучший среднеквадратичный прогноз \(X\) по
информации \(\mathcal G\):

\[
\mathbb E[(X-\mathbb E[X\mid\mathcal G])^2]
=
\min_{Y\in L^2(\mathcal G)}\mathbb E[(X-Y)^2].
\]

Главные ловушки: не путать ортогональность и независимость, не забывать
вторые моменты и замкнутость подпространства.

\subsection{15. Если осталось 2 минуты перед
экзаменом}\label{ux435ux441ux43bux438-ux43eux441ux442ux430ux43bux43eux441ux44c-2-ux43cux438ux43dux443ux442ux44b-ux43fux435ux440ux435ux434-ux44dux43aux437ux430ux43cux435ux43dux43eux43c}

Запомнить четыре формулы.

\textbf{\(L^2\) как гильбертово пространство:}

\[
\langle X,Y\rangle=\mathbb E[XY],
\qquad
\|X\|_2=(\mathbb EX^2)^{1/2}.
\]

\textbf{Ковариация как скалярное произведение:}

\[
\operatorname{Cov}(X,Y)
=
\langle X-\mathbb EX,\ Y-\mathbb EY\rangle.
\]

\textbf{Проектор:}

\[
P_Hx\in H,
\qquad
x-P_Hx\perp H.
\]

\textbf{Условное ожидание как проекция:}

\[
\mathbb E[X\mid\mathcal G]
=
P_{L^2(\mathcal G)}X.
\]

Главная ловушка: ортогональность/некоррелированность не означает
независимость; проекция гарантируется на замкнутое подпространство.

\newpage

\end{document}
